\section{Practice Activity 1}\label{sec:practice-activity-1}

Problems and detailed worked solutions on Expected Value, Maximin, and Maximax.

\subsection{Problem 1 --- Local Café Investment}\label{subsec:practice-activity-1-problem-1}

A small café is considering whether to introduce a new line of artisan desserts.
The owner has two options:

\begin{itemize}
	\item Launch the dessert line, in which case profits depend on demand:
	\begin{itemize}
		\item High demand \(\rightarrow\) 40 thousand USD
		\item Low demand \(\rightarrow\) 10 thousand USD
	\end{itemize}
	\item Keep the current menu, yielding 25 thousand USD regardless of demand.
\end{itemize}

Market research indicates:

\begin{itemize}
	\item Probability of high demand: 0.60
	\item Probability of low demand: 0.40
\end{itemize}

Question: Make a decision on the best option based on three perspectives: safe behavior (Maximin), balanced behavior (Expected Value), and risky behavior (Maximax).

\subsubsection{Solution 1 --- Local Café Investment}\label{subsec:practice-activity-1-solution-1}

\textbf{Structured solution}

\paragraph{(C2) Interpreting decision alternatives, events, consequences, and states}\label{subsubsec:practice-activity-1-solution-1-c2}

\begin{center}
	\begin{tabularx}{\linewidth}{@{}L{0.35\linewidth}Y@{}}
		\toprule
		Element & Description \\
		\midrule
		Alternatives & Launch desserts; Keep current menu \\
		States of nature & High demand (0.60); Low demand (0.40) \\
		Events & Actual demand level observed after the decision \\
		Consequences & Profit in thousands of USD for each alternative--state pair \\
		\bottomrule
	\end{tabularx}
\end{center}

\paragraph{(C3) Building the payoff table}\label{subsubsec:practice-activity-1-solution-1-c3}

\begin{center}
	\begin{tabularx}{\linewidth}{@{}L{0.3\linewidth}C C@{}}
		\toprule
		Alternative & High demand (0.60) & Low demand (0.40) \\
		\midrule
		Launch desserts & 40 & 10 \\
		Keep menu & 25 & 25 \\
		\bottomrule
	\end{tabularx}
\end{center}

\paragraph{(C4) Applying the Maximax criterion}\label{subsubsec:practice-activity-1-solution-1-c4}

The maximax rule ignores probabilities and selects the alternative with the highest possible payoff.

\[
\begin{aligned}
\max_i \text{Payoff}(\text{Launch}, S_i) &= \max\{40, 10\} = 40\\
\max_i \text{Payoff}(\text{Keep}, S_i) &= \max\{25, 25\} = 25
\end{aligned}
\]

\[
\max\{40, 25\} = 40
\]

Because 40 is the largest best payoff, the maximax (risk-seeking) recommendation is launch desserts.

\paragraph{(C5) Applying the Maximin criterion}\label{subsubsec:practice-activity-1-solution-1-c5}

The maximin rule focuses on the worst-case payoff of each alternative and picks the largest of those worst cases.

\[
\begin{aligned}
\min_i \text{Payoff}(\text{Launch}, S_i) &= \min\{40, 10\} = 10\\
\min_i \text{Payoff}(\text{Keep}, S_i) &= \min\{25, 25\} = 25
\end{aligned}
\]

\[
\max\{10, 25\} = 25
\]

Since 25 exceeds 10, the maximin (risk-averse) choice is keep the current menu.

\paragraph{(C1) Formulating the decision problem for an optimal choice}\label{subsubsec:practice-activity-1-solution-1-c1}

We frame the decision as choosing the menu strategy that maximizes the café's objective. Alternatives are
Launch desserts or Keep menu; the uncertain states are high or low demand with
probabilities 0.60 and 0.40. The optimal decision under each criterion uses the following calculations.

\[
\begin{aligned}
EV_{\text{Launch}} &= 0.60(40) + 0.40(10) \\
&= 0.60\cdot40 + 0.40\cdot10 \\
&= 24 + 4 = 28 \\
EV_{\text{Keep}} &= 0.60(25) + 0.40(25) \\
&= 0.60\cdot25 + 0.40\cdot25 \\
&= 15 + 10 = 25
\end{aligned}
\]

\[
\max\{28, 25\} = 28
\]

Summary of recommendations across decision criteria

\begin{itemize}
	\item Maximax favors launch desserts (highest payoff of 40).
	\item Maximin favors keep the current menu (best worst-case payoff of 25).
	\item Expected value favors launch desserts (EV of 28).
\end{itemize}

Comparing the criteria shows that the launch option leads on optimistic and average performance, while keeping the
menu provides stronger downside protection.

\subsection{Problem 2 --- Green Energy Production Mix}\label{subsec:practice-activity-1-problem-2}

A renewable-energy company is evaluating its energy production strategy for the coming year.
Three weather conditions are possible:

\begin{itemize}
	\item Windy
	\item Sunny
	\item Cloudy
\end{itemize}

The company is comparing three strategies:

\begin{enumerate}
	\item Build wind turbines
	\item Build solar farms
	\item Build a balanced hybrid system
\end{enumerate}

Expected profits (in million USD):

\begin{itemize}
	\item Wind turbines \(\rightarrow\) 85 (Windy), 40 (Sunny), 25 (Cloudy)
	\item Solar farms \(\rightarrow\) 30 (Windy), 90 (Sunny), 20 (Cloudy)
	\item Hybrid system \(\rightarrow\) 60 (Windy), 60 (Sunny), 55 (Cloudy)
\end{itemize}

Weather probabilities:

\begin{itemize}
	\item Windy \(\rightarrow\) 0.30
	\item Sunny \(\rightarrow\) 0.50
	\item Cloudy \(\rightarrow\) 0.20
\end{itemize}

Question: Make a decision on the best strategy based on three perspectives: safe behavior (Maximin), balanced behavior (Expected Value), and risky behavior (Maximax).

\subsubsection{Solution 2 --- Green Energy Production Mix}\label{subsec:practice-activity-1-solution-2}

\textbf{Structured solution}

\paragraph{(C2) Interpreting decision alternatives, events, consequences, and states}\label{subsubsec:practice-activity-1-solution-2-c2}

\begin{center}
	\begin{tabularx}{\linewidth}{@{}L{0.35\linewidth}Y@{}}
		\toprule
		Element & Description \\
		\midrule
		Alternatives & Wind turbines; Solar farms; Hybrid system \\
		States of nature & Windy (0.30); Sunny (0.50); Cloudy (0.20) \\
		Events & Weather outcome driving power generation \\
		Consequences & Profit in million USD for each strategy--weather combination \\
		\bottomrule
	\end{tabularx}
\end{center}

\paragraph{(C3) Building the payoff table}\label{subsubsec:practice-activity-1-solution-2-c3}

\begin{center}
	\begin{tabularx}{\linewidth}{@{}L{0.3\linewidth}C C C@{}}
		\toprule
		Strategy & Windy (0.30) & Sunny (0.50) & Cloudy (0.20) \\
		\midrule
		Wind turbines & 85 & 40 & 25 \\
		Solar farms & 30 & 90 & 20 \\
		Hybrid system & 60 & 60 & 55 \\
		\bottomrule
	\end{tabularx}
\end{center}

\paragraph{(C4) Applying the Maximax criterion}\label{subsubsec:practice-activity-1-solution-2-c4}

Maximax considers only the best possible outcome for each alternative and then picks the largest one.

\[
\begin{aligned}
\max_i \text{Payoff}(\text{Wind}, S_i) &= \max\{85, 40, 25\} = 85\\
\max_i \text{Payoff}(\text{Solar}, S_i) &= \max\{30, 90, 20\} = 90\\
\max_i \text{Payoff}(\text{Hybrid}, S_i) &= \max\{60, 60, 55\} = 60
\end{aligned}
\]

\[
\max\{85, 90, 60\} = 90
\]

The largest best payoff is 90, so the maximax (risk-seeking) decision is build solar farms.

\paragraph{(C5) Applying the Maximin criterion}\label{subsubsec:practice-activity-1-solution-2-c5}

Maximin evaluates the worst payoff for each alternative and favors the highest of those worst cases.

\[
\begin{aligned}
\min_i \text{Payoff}(\text{Wind}, S_i) &= \min\{85, 40, 25\} = 25\\
\min_i \text{Payoff}(\text{Solar}, S_i) &= \min\{30, 90, 20\} = 20\\
\min_i \text{Payoff}(\text{Hybrid}, S_i) &= \min\{60, 60, 55\} = 55
\end{aligned}
\]

\[
\max\{25, 20, 55\} = 55
\]

Because 55 is the largest minimum payoff, the maximin (risk-averse) recommendation is the hybrid system.

\paragraph{(C1) Formulating the decision problem for an optimal choice}\label{subsubsec:practice-activity-1-solution-2-c1}

The firm wants the production mix that best supports its profit goal. Alternatives are wind,
solar, or a hybrid system and the uncertain states are windy,
sunny, or cloudy with respective probabilities 0.30, 0.50, and 0.20. Optimizing
by expected value yields:

\[
\begin{aligned}
EV_{\text{Wind}} &= 0.30(85) + 0.50(40) + 0.20(25) \\
&= 0.30\cdot85 + 0.50\cdot40 + 0.20\cdot25 \\
&= 25.5 + 20 + 5 = 50.5 \\
EV_{\text{Solar}} &= 0.30(30) + 0.50(90) + 0.20(20) \\
&= 0.30\cdot30 + 0.50\cdot90 + 0.20\cdot20 \\
&= 9 + 45 + 4 = 58 \\
EV_{\text{Hybrid}} &= 0.30(60) + 0.50(60) + 0.20(55) \\
&= 0.30\cdot60 + 0.50\cdot60 + 0.20\cdot55 \\
&= 18 + 30 + 11 = 59
\end{aligned}
\]

\[
\max\{50.5, 58, 59\} = 59
\]

Summary of recommendations across decision criteria

\begin{itemize}
	\item Maximax favors solar farms (highest payoff of 90).
	\item Maximin favors the hybrid system (best worst-case payoff of 55).
	\item Expected value favors the hybrid system (EV of 59).
\end{itemize}

The comparison shows solar excels only in the best-case outcome, while the hybrid strategy dominates on downside
protection and average profitability.

\subsection{Problem 3 --- Global Shipping Network Design}\label{subsec:practice-activity-1-problem-3}

A multinational logistics firm must redesign its global shipping network for the next five years to balance efficiency
with resilience. Management is weighing three mutually exclusive configurations: a centralized network built around a
single mega-hub in East Asia, a regional network with three hubs in the Americas, Europe, and Asia, and an adaptive
AI--driven network that relies on dynamic routing powered by machine-learning models.

The firm is uncertain about future trade volatility. It models three possible scenarios: low volatility with stable trade
flows, medium volatility with moderate regional disruptions, and high volatility with significant unpredictability. The
firm assigns probabilities of 0.25 to low volatility, 0.45 to medium volatility, and 0.30 to high volatility. Expected
profits (in million USD) are estimated at 210, 130, and $-40$ for the centralized option under low, medium, and high
volatility; 180, 160, and 70 for the regional option; and 150, 175, and 140 for the AI--driven option.

Question: Make a decision on the best configuration based on three perspectives: safe behavior (Maximin), balanced behavior (Expected Value), and risky behavior (Maximax).

\subsubsection{Solution 3 --- Global Shipping Network Design}\label{subsec:practice-activity-1-solution-3}

\textbf{Structured solution}

\paragraph{(C2) Interpreting decision alternatives, events, consequences, and states}\label{subsubsec:practice-activity-1-solution-3-c2}

\begin{center}
	\begin{tabularx}{\linewidth}{@{}L{0.35\linewidth}Y@{}}
		\toprule
		Element & Description \\
		\midrule
		Alternatives & Centralized hub; Regional hubs; AI--driven adaptive network \\
		States of nature & Low volatility (0.25); Medium (0.45); High (0.30) \\
		Events & Actual trade volatility realized over the planning horizon \\
		Consequences & Profit in million USD for each configuration--volatility scenario \\
		\bottomrule
	\end{tabularx}
\end{center}

\paragraph{(C3) Building the payoff table}\label{subsubsec:practice-activity-1-solution-3-c3}

\begin{center}
	\begin{tabularx}{\linewidth}{@{}L{0.3\linewidth}C C C@{}}
		\toprule
		Configuration & Low volatility (0.25) & Medium volatility (0.45) & High volatility (0.30) \\
		\midrule
		Centralized & 210 & 130 & -40 \\
		Regional & 180 & 160 & 70 \\
		AI--Driven & 150 & 175 & 140 \\
		\bottomrule
	\end{tabularx}
\end{center}

\paragraph{(C4) Applying the Maximax criterion}\label{subsubsec:practice-activity-1-solution-3-c4}

Maximax selects the alternative with the highest possible payoff regardless of downside.

\[
\begin{aligned}
\max_i \text{Payoff}(\text{Centralized}, S_i) &= \max\{210, 130, -40\} = 210\\
\max_i \text{Payoff}(\text{Regional}, S_i) &= \max\{180, 160, 70\} = 180\\
\max_i \text{Payoff}(\text{AI}, S_i) &= \max\{150, 175, 140\} = 175
\end{aligned}
\]

\[
\max\{210, 180, 175\} = 210
\]

The greatest best payoff is 210, so the maximax choice is the centralized network.

\paragraph{(C5) Applying the Maximin criterion}\label{subsubsec:practice-activity-1-solution-3-c5}

Maximin emphasizes protection against the worst case by comparing the minimum payoff of each alternative.

\[
\begin{aligned}
\min_i \text{Payoff}(\text{Centralized}, S_i) &= \min\{210, 130, -40\} = -40\\
\min_i \text{Payoff}(\text{Regional}, S_i) &= \min\{180, 160, 70\} = 70\\
\min_i \text{Payoff}(\text{AI}, S_i) &= \min\{150, 175, 140\} = 140
\end{aligned}
\]

\[
\max\{-40, 70, 140\} = 140
\]

Because 140 is the largest of these minimums, the maximin recommendation is the AI--driven network.

\paragraph{(C1) Formulating the decision problem for an optimal choice}\label{subsubsec:practice-activity-1-solution-3-c1}

The logistics firm frames its decision as selecting the network that maximizes long-run profit. Alternatives are
centralized, regional, or AI--driven and the uncertain states are
low, medium, or high volatility with probabilities 0.25, 0.45, and
0.30. Using expected value to target optimality:

\[
\begin{aligned}
EV_{\text{Centralized}} &= 0.25(210) + 0.45(130) + 0.30(-40) \\
&= 0.25\cdot210 + 0.45\cdot130 + 0.30\cdot(-40) \\
&= 52.5 + 58.5 - 12 = 99 \\
EV_{\text{Regional}} &= 0.25(180) + 0.45(160) + 0.30(70) \\
&= 0.25\cdot180 + 0.45\cdot160 + 0.30\cdot70 \\
&= 45 + 72 + 21 = 138 \\
EV_{\text{AI}} &= 0.25(150) + 0.45(175) + 0.30(140) \\
&= 0.25\cdot150 + 0.45\cdot175 + 0.30\cdot140 \\
&= 37.5 + 78.75 + 42 = 158.25
\end{aligned}
\]

\[
\max\{99, 138, 158.25\} = 158.25
\]

Summary of recommendations across decision criteria

\begin{itemize}
	\item Maximax favors the centralized network (highest payoff of 210).
	\item Maximin favors the AI--driven network (best worst-case payoff of 140).
	\item Expected value favors the AI--driven network (EV of 158.25).
\end{itemize}

The centralized design wins only on peak upside, while the AI--driven option provides the strongest performance on
both downside risk and average returns across volatility scenarios.

\vspace{6pt}
\noindent\hyperref[table-of-contents]{Back to Top}
